\chapter{流式备份：512KB 一段段读}
\label{ch:gen1-streaming}

\section{为什么不能一次读完整文件？}

大文件可能有几个 GB。一次性读入内存不现实。  
\textbf{流式 CDC}：每次只读约 \textbf{512KB}（叫一个 feed），在这 512KB 里继续滚指纹、找切点。

\begin{figure}[htbp]
\centering
\begin{tikzpicture}[node distance=0.15cm]
  \node[note] at (0,2.1) {整文件（例如 2GB）——逻辑上是一长条字节流};

  \node[feed, minimum width=2.4cm] (f0) at (0,1) {feed 0\\512KB};
  \node[feedalt, minimum width=2.4cm, right=0.12cm of f0] (f1) {feed 1\\512KB};
  \node[feed, minimum width=2.4cm, right=0.12cm of f1] (f2) {feed 2\\512KB};
  \node[feedalt, minimum width=2.4cm, right=0.12cm of f2] (f3) {feed 3\\512KB};
  \node[right=0.12cm of f3, font=\Large] {$\cdots$};

  \draw[cutline] ($(f1.east)!0.55!(f2.west)$) -- ++(0,1.35);
  \node[font=\scriptsize, CutRed, align=center] at ($(f1.east)!0.55!(f2.west)+(0,1.55)$)
    {切点可以落在\\feed 1 与 2 之间};

  \node[zone=Gen1Blue, minimum width=10.5cm, below=0.85cm of f1]
    (note) {对算法而言仍是\textbf{连续}字节流；feed 只是 IO 读入的单位};
\end{tikzpicture}
\caption{切点可以跨 feed；所以必须保存滚动指纹 $h$ 和 carry 尾巴。}
\end{figure}

\section{Carry：跨段的「尾巴」}

如果一个 chunk 还没切完，512KB 读完了怎么办？  
把还没决定的部分（tail）存进 \keyterm{carry}，和下一段 feed 拼在一起继续扫。

\begin{figure}[htbp]
\centering
\begin{tikzpicture}[node distance=0.35cm]
  \node[gen1, minimum width=2.8cm] (f0) {feed 0\\512KB};
  \node[gen1, minimum width=2.8cm, right=0.5cm of f0] (carry) {carry\\尾巴};
  \node[gen1, minimum width=2.8cm, right=0.5cm of carry] (f1) {feed 1\\512KB};
  \draw[arr] (f0) -- node[above,font=\tiny]{未切完} (carry);
  \draw[arr] (carry) -- node[above,font=\tiny]{拼接} (f1);
\end{tikzpicture}
\caption{carry + 新 feed = 逻辑上连续的字节流（对算法而言）。}
\end{figure}

\section{流式状态机（第一代）}

\begin{figure}[htbp]
\centering
\begin{tikzpicture}[node distance=1.1cm]
  \node[gen1, minimum width=2.4cm] (read) {读 feed};
  \node[gen1, minimum width=2.6cm, right=of read] (merge) {carry 拼接};
  \node[gen1, minimum width=2.6cm, right=of merge] (scan) {CDC 扫描};
  \node[gen2, minimum width=2.4cm, right=of scan] (out) {输出 chunk};
  \node[gen1, minimum width=2.4cm, below=0.9cm of scan] (save) {保存 resume\\$h$ + carry};

  \draw[arr] (read) -- (merge);
  \draw[arr] (merge) -- (scan);
  \draw[arr] (scan) -- (out);
  \draw[arr] (scan) -- (save);
  \draw[arr] (save.west) -| (read.south);
\end{tikzpicture}
\caption{每处理完一个 feed，把「没切完的状态」写进 resume，下一轮接着滚。}
\end{figure}

\section{Resume 状态（第一代）}

流式程序必须记住：

\begin{itemize}[nosep]
  \item 当前滚动指纹 $h$（若扫描中途停下）；
  \item carry 里有哪些字节；
  \item 文件绝对位置（第几个字节了）。
\end{itemize}

\begin{figure}[htbp]
\centering
\begin{tikzpicture}[node distance=0.4cm]
  \node[gen1, minimum width=3.2cm, align=left] (r1) {resume.h\\{\footnotesize 32 位指纹}};
  \node[gen1, minimum width=3.2cm, align=left, right=0.5cm of r1] (r2)
    {resume.carry\\{\footnotesize 字节数组}};
  \node[gen1, minimum width=3.2cm, align=left, right=0.5cm of r2] (r3)
    {resume.pos\\{\footnotesize 绝对偏移}};
  \node[below=0.6cm of r2, font=\scriptsize, align=center] (cap)
    {三者齐全 → 下一 feed 接上后切点与整文件 CDC 一致};
\end{tikzpicture}
\caption{resume 就是流式 CDC 的「存档点」。}
\end{figure}

\begin{stepbox}[Step 6 — 黄金规则]
\textbf{整文件一次性 CDC} 和 \textbf{512KB 流式 CDC} 必须切在\textbf{完全相同}的位置。  
这叫 \keyterm{parity}（一致性）。第二代也要遵守。
\end{stepbox}

\section{和 pipeline 的关系}

备份时往往：一边 CDC，一边算哈希、一边压缩。  
第一代 Stream 路径已经高度优化；第二代 2F-Gear 要接在同一套流式框架里（只是多记一个指纹）。

\begin{figure}[htbp]
\centering
\begin{tikzpicture}[node distance=0.45cm]
  \node[byte, minimum width=1.6cm] (in) {文件};
  \node[gen1, right=0.5cm of in, minimum width=1.8cm] (cdc) {CDC\\切块};
  \node[gen1, right=0.5cm of cdc, minimum width=1.8cm] (hash) {哈希};
  \node[gen1, right=0.5cm of hash, minimum width=1.8cm] (cmp) {压缩};
  \node[gen2, right=0.5cm of cmp, minimum width=1.8cm] (store) {入库};
  \draw[arr] (in) -- (cdc) -- (hash) -- (cmp) -- (store);
  \node[below=0.55cm of cdc, font=\scriptsize, align=center] (note)
    {本手册只讲这一块\\（Stream 或 2F-Gear）};
\end{tikzpicture}
\caption{CDC 在最前面：切对了，后面 dedup 和压缩才有意义。}
\end{figure}

\section{模拟示例：mini-feed（每段 4 字节）}

\begin{simbox}[12 字节文件分成 3 个 feed]
\begin{tabular}{@{}cll@{}}
\toprule
\textbf{feed} & \textbf{字节} & \textbf{事件} \\
\midrule
0 & 0--3 & 扫到 $i=3$ 切点 → 输出 chunk1 \\
1 & 4--7 & 新 chunk；在 $i=7$ 切 \\
2 & 8--11 & 若未切完 → carry 拼 feed 3 \\
\bottomrule
\end{tabular}

\vspace{0.4em}
真实备份 feed=512KB，逻辑相同：\textbf{切点跨 feed 时 resume 必须存 $h$}。
\end{simbox}
